\documentclass[10pt,aspectratio=169]{beamer}

% ----------------------- Packages -----------------------
\usepackage[english]{babel}
\usepackage[utf8]{inputenc}
\usepackage{graphicx}
\usepackage{amsmath,amssymb}
\usepackage{pgfplots}
\usepackage{pgfplotstable}
\pgfplotsset{compat=1.18}
\usepackage{multicol}
\usepackage{subfigure}
\usepackage{listings}
\usepackage[UKenglish]{isodate}
\usepackage{tikz}
\usetikzlibrary{arrows.meta,positioning,fit,shapes.geometric,calc,fillbetween}

\cleanlookdateon

% ----------------------- Listings Configuration -----------------------
\lstset{
  basicstyle=\ttfamily\small\color{white},
  backgroundcolor=\color{black!80},
  keywordstyle=\color{cyan!80}\bfseries,
  commentstyle=\color{white!60}\itshape,
  stringstyle=\color{green!60},
  numberstyle=\tiny\color{white!40},
  numbers=left,
  numbersep=8pt,
  frame=single,
  rulecolor=\color{white!30},
  breaklines=true,
  breakatwhitespace=true,
  tabsize=2,
  showstringspaces=false
}

% ----------------------- Theme & Colors -----------------------
\usetheme{Madrid}
\usecolortheme{owl}   % make sure you have the beamer-owl theme installed

% Dark background + white text
\setbeamercolor{background canvas}{bg=black}
\setbeamercolor{normal text}{fg=white}
\setbeamercolor{frametitle}{fg=white,bg=black}
\setbeamercolor{title}{fg=white}
\setbeamercolor{author}{fg=white}
\setbeamercolor{institute}{fg=white}
\setbeamercolor{date}{fg=white}

% ----------------------- Metadata -----------------------
\title[Short Title]{Main Talk Title Goes Here}
\author{Presenter Name}
\institute{Your Institution}
\date{\today} % or a fixed date

% ----------------------- Left Sidebar Progress Rail -----------------------
\useoutertheme[left,width=18mm]{sidebar}

% Fonts in the rail
\setbeamerfont{title in sidebar}{size=\tiny}
\setbeamerfont{author in sidebar}{size=\tiny}
\setbeamerfont{section in sidebar}{size=\scriptsize}
\setbeamerfont{subsection in sidebar}{size=\tiny}

% Colors in the rail
\setbeamercolor{sidebar}{bg=black, fg=white}
\setbeamercolor{title in sidebar}{fg=white}
\setbeamercolor{section in sidebar}{fg=white,bg=black}
\setbeamercolor{section in sidebar shaded}{fg=white!40,bg=black}
\setbeamercolor{subsection in sidebar}{fg=white}
\setbeamercolor{subsection in sidebar shaded}{fg=white!35}
\setbeamercolor{section in sidebar current}{fg=black,bg=cyan!60}
\setbeamercolor{subsection in sidebar current}{fg=black,bg=cyan!60}

\makeatletter
\setbeamertemplate{sidebar canvas left}{%
  \color{black}\vrule width\beamer@sidebarwidth height\paperheight
}

\setbeamertemplate{sidebar left}{%
  \vskip13ex%
  \hbox to \beamer@sidebarwidth{%
    \hfil%
    \vbox to\paperheight{%
      \begin{minipage}[t]{\beamer@sidebarwidth}%
        \vspace*{1ex}
        \raggedright
        \insertverticalnavigation{\paperheight}% sections/subsections
        \vfill
      \end{minipage}%
    }%
    \hfil%
  }%
}
\makeatother

% Remove default headline
\setbeamertemplate{headline}{}

% Custom frametitle bar across full width
\setbeamertemplate{frametitle}{%
  \vskip0.5ex
  \begin{beamercolorbox}[wd=\paperwidth,ht=2.8ex,dp=1ex,leftskip=1em]{frametitle}
    \usebeamerfont{frametitle}\insertframetitle
  \end{beamercolorbox}
}
\setbeamercolor{frametitle}{fg=white,bg=black}

% ----------------------- Document -----------------------
\begin{document}

% ----------------------- Title frame (shifted to center, ignoring sidebar) -----------------------
\begin{frame}[plain]
  \begingroup
    % shift content left by the known sidebar width you set: width=18mm
    \hspace*{-25mm}%
    \begin{minipage}{\paperwidth}
      \centering
      \titlepage
    \end{minipage}
  \endgroup
\end{frame}

% =======================================================
% SECTION 1: INTRODUCTION
% =======================================================
\section{Introduction}

\subsection{Motivation}
% ---------- Simple bullet slide ----------
\begin{frame}{Motivation}
  \begin{itemize}
    \item Capture the audience's attention, since this is all you need \cite{vaswani2023attentionneed}.
    \item Citations are supported as seen in the line above.
    \item Briefly explain why this topic matters.
    \item State the real-world problem or use case.
    \item Optionally mention who might benefit from your results.
  \end{itemize}
\end{frame}


% ---------- Example: Figure slide ----------
\subsection{Figure}
\begin{frame}{Figure}
  \begin{figure}
    \centering
    % Replace `example-image` with your own image file
    \includegraphics[width=0.5\textwidth]{example-image}
    \caption{Short descriptive caption of your figure.}
  \end{figure}
\end{frame}

% =======================================================
% SECTION 2: BACKGROUND / THEORY
% =======================================================
\section{Background}

\subsection{Key Concepts}
% ---------- Example: Two-column slide ----------
\begin{frame}{Key Concepts}
  \begin{columns}[T,totalwidth=\textwidth]
    \begin{column}{0.48\textwidth}
      \textbf{Concept A}
      \begin{itemize}
        \item Short definition.
        \item Why it is important.
      \end{itemize}
    \end{column}
    \begin{column}{0.48\textwidth}
      \textbf{Concept B}
      \begin{itemize}
        \item Short definition.
        \item Relation to Concept A.
      \end{itemize}
    \end{column}
  \end{columns}
\end{frame}

% ---------- Example: Simple TikZ diagram ----------
\begin{frame}{Example TikZ Diagram}
  \centering
  \begin{tikzpicture}[
    >=Stealth,
    node distance=2cm,
    box/.style={rectangle, rounded corners, draw=white, thick,
                fill=black!55, minimum width=2.6cm,
                minimum height=0.9cm, text=white, align=center},
    arrow/.style={->, thick, white}
  ]
    \node[box] (input)  {Input};
    \node[box, right=of input] (process) {Process};
    \node[box, right=of process] (output) {Output};

    \draw[arrow] (input) -- (process);
    \draw[arrow] (process) -- (output);
  \end{tikzpicture}

  \vspace{1em}
  \begin{itemize}
    \item Use boxes and arrows to show the logical flow.
    \item Replace labels with your actual pipeline stages.
  \end{itemize}
\end{frame}

% ---------- Example: Code listing ----------
\subsection{Code Example}
\begin{frame}[fragile]{Code Listing Example}
  \begin{lstlisting}[language=Python, basicstyle=\ttfamily\footnotesize\color{white}]
    def fibonacci(n):
        """Calculate the nth Fibonacci number."""
        if n <= 1:
            return n
        else:
            return fibonacci(n-1) + fibonacci(n-2)

    # Example usage
    result = fibonacci(10)
    print(f"The 10th Fibonacci number is {result}")
  \end{lstlisting}
  \vspace{0.5em}
  \small{Replace with your own code. Supports many languages: Python, Java, C++, etc.}
\end{frame}

% ---------- Example: Block environments ----------
\subsection{Block Environments}
\begin{frame}{Block Environments}
  \begin{block}{Standard Block}
    Use blocks to highlight important information or definitions.
  \end{block}

  \vspace{0.3em}

  \begin{alertblock}{Alert Block}
    Use alert blocks for warnings, cautions, or critical points.
  \end{alertblock}

  \vspace{0.3em}

  \begin{exampleblock}{Example Block}
    Use example blocks to show concrete examples or use cases.
  \end{exampleblock}

  \vspace{0.5em}
  \small{These colored blocks help organize and emphasize content.}
\end{frame}

% =======================================================
% SECTION 3: METHOD / EXPERIMENTS
% =======================================================
\section{Methodology}

\subsection{Overview}
% ---------- Example: Method bullet slide ----------
\begin{frame}{Method Overview}
  \begin{enumerate}
    \item \textbf{Data / Setting:} Describe the dataset or experimental setup.
    \item \textbf{Model / Algorithm:} Summarize what you apply.
    \item \textbf{Metrics:} Specify how you evaluate performance.
  \end{enumerate}
\end{frame}

\subsection{Some Plot}
% ---------- Example: PGFPlots chart ----------
\begin{frame}{Example PGFPlots Chart}
  % Small dummy table inline
  \pgfplotstableread[row sep=\\,col sep=&]{
  x   & y   \\
  1   & 2   \\
  2   & 3.5 \\
  3   & 3.2 \\
  4   & 4.8 \\
  }\datatable

  \begin{tikzpicture}
    \begin{axis}[
      width=0.9\textwidth,
      height=0.45\textheight,
      xlabel={X axis label},
      ylabel={Y axis label},
      axis line style={white},
      tick align=inside,
      tick style={white},
      xticklabel style={text=white},
      yticklabel style={text=white},
      grid=major,
      grid style={white!20},
      label style={text=white}
    ]
      \addplot+[mark=o, thick] table[x=x,y=y]{\datatable};
    \end{axis}
  \end{tikzpicture}

  \vspace{0.5em}
  \small{Replace the data table and labels with your own results.}
\end{frame}

% =======================================================
% SECTION 4: RESULTS & DISCUSSION
% =======================================================
\section{Results}

\subsection{Key Findings}
% ---------- Example: Result bullets ----------
\begin{frame}{Key Results}
  \begin{itemize}
    \item \textbf{Result 1:} Short statement of the main finding.
    \item \textbf{Result 2:} Compare against baseline or prior work.
    \item \textbf{Result 3:} Mention surprising or counterintuitive observations.
  \end{itemize}
\end{frame}

% ---------- Example: Results table ----------
\subsection{Comparison Table}
\begin{frame}{Performance Comparison}
  \begin{table}
    \centering
    \begin{tabular}{lccc}
      \hline
      \textbf{Method} & \textbf{Accuracy} & \textbf{Time (s)} & \textbf{Memory (MB)} \\
      \hline
      Baseline        & 85.3\%            & 12.5              & 256 \\
      Method A        & 89.7\%            & 10.2              & 312 \\
      Method B        & 92.1\%            & 15.8              & 289 \\
      \textbf{Ours}   & \textbf{94.5\%}   & \textbf{9.8}      & \textbf{278} \\
      \hline
    \end{tabular}
    \caption{Quantitative comparison of different methods.}
  \end{table}

  \vspace{0.5em}
  \small{Our method achieves the best accuracy with competitive speed and memory usage.}
\end{frame}

\section{Discussion}

\begin{frame}{Discussion}
  \begin{itemize}
    \item Interpret what the results mean.
    \item Limitations of your approach or dataset.
    \item Possible future directions or open questions.
  \end{itemize}
\end{frame}

% =======================================================
% SECTION 5: CONCLUSION
% =======================================================
\section{Conclusion}

\begin{frame}{Conclusion}
  \begin{itemize}
    \item Summarize the core message of the talk in 2–3 bullets.
    \item Highlight the main contribution or takeaway.
    \item Optionally include a call to action or next steps.
  \end{itemize}
\end{frame}

% ---------- Thank you / Questions ----------
\begin{frame}
  \centering
  \Large Thank you for your attention!\\[1em]
  \large Questions?
\end{frame}

% =======================================================
% REFERENCES (optional – uncomment if you use BibTeX)
% =======================================================
\section{References}
\begin{frame}[allowframebreaks]{References}
  \bibliographystyle{splncs04}
  % add your .bib file here
  \bibliography{lit}
\end{frame}

\end{document}
